\section{News explains little of a stock's own move; the commercial models are overconfident}
\label{sec:explaining}

\subsection{Explaining realized moves}\label{sec:ladder}

\citet{roll1988} asked how much of a stock's daily return public news explains, and \citet{kelly2026} revisit
the question for the market with generative AI. We ask it for single stocks, out of sample. For every covered
company-day of the test period (days without selected news enter with zero news features, as in Roll), we build
up the explanation of the stock-specific move $z_{cc}$ step by step: the market and the sector (Roll's first
step), then the pre-open news reading, then the reading of all the day's news, and finally the post-reaction
reading's explanation of the move, which reads the realized move and is therefore an attribution, not a
forecast. Each step is a ridge regression fitted out of sample by the pre-registered protocol (validation data
for the fit and the penalty), and each is scored on the test by
\begin{equation}\label{eq:r2}
R^{2}_{\mathrm{OOS}} = 1 - \frac{\sum_{i,t}\bigl(y_{i,t} - \hat y_{i,t}\bigr)^{2}}
{\sum_{i,t}\bigl(y_{i,t} - \bar y_{\mathrm{ref}}\bigr)^{2}},
\qquad
\Delta R^{2}_{s} = R^{2}_{\mathrm{OOS}}(s) - R^{2}_{\mathrm{OOS}}(s-1),
\end{equation}
where $\bar y_{\mathrm{ref}}$ is the reference the frozen specification fixes (a zero forecast for the
stock-specific move, and the val-B mean for the size of the move in H4), $R^{2}_{\mathrm{OOS}}(s)$ is the fit of
steps one to $s$, and $\Delta R^{2}_{s}$ is what step $s$ adds to the steps before it. A positive $R^{2}_{\mathrm{OOS}}$ means the step predicts better than the reference
on data it was not fitted on.

The news explains little of a stock's own move (Figure~\ref{fig:explaining}, Table~\ref{tab:explaining}). The
market and the sector explain \res{roll-r2-pct} of the total daily variance. Of the stock-specific part, the
pre-open reading explains \res{lf-ladder-ante-pct} for Laya-Fin and \res{d1-ladder-ante-pct} for d1-ultra, and
the reading of all the day's news \res{lf-ladder-post-pct} and \res{d1-ladder-post-pct}. H5, the secondary
hypothesis that the same-day news step is positive, is supported for Laya-Fin
($q = \res{lf-h5-q}$), the only secondary test that survives Benjamini--Hochberg for it, and for d1-ultra
($q \res{d1-h5-q}$). Much of what the news explains was already knowable before the open
(\res{lf-knowable-pct} for Laya-Fin, \res{d1-knowable-pct} for d1-ultra). Reading the realized move explains far more: d1-ultra's
explanation step reaches \res{d1-ladder-expl-pct} of the stock-specific variance on all covered days and
\res{d1-ladder-expl-event-pct} on days with news.

\begin{figure}[tbp]
\centering
\includegraphics{\figdir/fig5_explaining_moves.pdf}
\caption{\textbf{Same-day news explains \res{lf-ladder-post-pct} of the stock-specific daily variance out of
sample for Laya-Fin and \res{d1-ladder-post-pct} for d1-ultra; reading the realized move explains
\res{d1-ladder-expl-pct}.} Test period; out-of-sample $R^2$ in percent of the stock-specific variance (the market
and the sector removed). (a) d1-ultra's $R^2$ by cut, stacked by step in the legend's order: the pre-open news,
plus the day's news, plus the post-reaction explanation (which reads the realized move, so it is not a
forecast). Each total stands above its bar with its 95\% bootstrap interval; the two thin news steps are
labelled beside each bar. For comparison, the market and the sector explain \res{roll-r2-pct} of the total
daily variance, a different base. (b) The news steps on their own scale, with 95\% bootstrap intervals, for
Laya-Fin and d1-ultra (model-specific). Only d1-ultra answers the explanation question; the anchor's near-constant
readings are in Table~\ref{tab:explaining}.}\label{fig:explaining}
\end{figure}

\input{\tabdir/tab5_explaining_moves.tex}

Being explained does not predict persistence. Days the model judges explained (the probability that the move
is explained by the news is at least one half) are \res{expl-events-share} of event days but carry
\res{expl-var-share} of their stock-specific variance. Yet in the pre-registered persistence regression the
coefficient on the explanation is \res{theta-d2} after the close (standard error \res{theta-d2-se},
$p = \res{theta-d2-p}$ for a positive effect): explained moves do not persist more than unexplained ones, and if
anything reverse slightly more. The size of the move is easier to read than its direction, as
Section~\ref{sec:validation} showed for H4: the models rank which stocks will move a lot, but add nothing to
recent volatility.

\subsection{Calibration}\label{sec:calibration}

The pre-registered calibration metric is the expected calibration error \citep{pakdamannaeini2015,guo2017}: the
pre-open reading's nine levels are grouped into three classes (a fall, no move and a rise, split at half a
typical daily move), and the top class's stated probability is compared with its hit rate against the realized
five-day move, in 15 equal-width bins. Laya-Fin is calibrated but compressed: its error is \res{lf-ece} and its
log-loss improves on the validation base rates by only \res{lf-llskill}, because its probabilities stay close to
those base rates. The commercial models are overconfident: their stated probabilities overstate their hit rates,
with errors of \res{d1-ece} (d1-ultra), \res{span-ece} (span-02) and \res{drex-ece} (drex-v2.0) and negative
log-loss skill. A single temperature fitted on val-A lowers d1-ultra's error from \res{d1-ece-before} to
\res{d1-ece-after} (Appendix Table~\ref{tab:a-calibration} and Figure~\ref{fig:a-reliability}). Fine-tuning
does not buy a significant IC gain over the untuned model either: H7's difference between Laya-Fin and Laya
zero-shot is \res{lf-h7} ($p = \res{lf-h7-p}$, $q = \res{lf-h7-q}$). The cost per thousand decisions is zero for
the trained models on a free GPU tier, and \res{span-cost} to \res{d1-cost} dollars for the commercial models.
